\documentclass[../../../include/open-logic-section]{subfiles}

\begin{document}

\olfileid{sth}{spine}{recursion}
\olsection{لامتناہی درجیاں والی بازگشت دے قضیے}

پر سانوں سبھ توں پہلاں ایہ پرکھنا اے کہ \olref[valpha]{defValphas}
اک کامیاب تعریف اے۔ ہور عام طور اُتے سانوں ثابت کرنا اے کہ لامتناہی
درجیاں والی بازگشت نال تعریف دین دی ہر کوشش کامیاب ہووے گی۔ ایس حصے دا ایہو مقصد اے۔

\emph{دھیان رہے۔ ایہ کم اوکھا اے}۔ پر ساری گل دا حاصل بڑا سدھا اے۔
لامتناہی درجیاں اُتے استقرا تے بدل رکھن دا اصول، لامتناہی درجیاں والی
بازگشت دے (کئی روپاں دے) جائز ہون دی ضمانت دیندے نیں۔\footnote{یاد رہے
کہ سارے فارمولیاں تے ٹرماں دے پیرامیٹر ہو سکدے نیں (جدوں تک صاف طور اُتے
کجھ ہور نہ آکھیا جاوے)۔}

\begin{defn}
منّ لو $\tau(x)$ اک ٹرم اے، $f$ اک فنکشن اے، تے $\alpha$ اک ترتیبی عدد اے۔
اسیں $f$ نوں \emph{$\alpha$ تک دا لگ بھگ روپ} آکھدے آں، $\tau$ لئی
اُدوں تے صرف اُدوں جدوں $\dom{f} = \alpha$ تے
$(\forall \beta \in \alpha)f(\beta) = \tau(\funrestrictionto{f}{\beta})$
دونیں شرطاں پوریاں ہون۔
\end{defn}

\begin{lem}[حدّ بندی والی بازگشت]\ollabel{transrecursionfun}
کسے وی ٹرم $\tau(x)$ تے کسے وی ترتیبی عدد $\alpha$ لئی، $\alpha$ تک دا
اکّو اک لگ بھگ روپ اے، $\tau$ لئی۔
\end{lem}
\begin{proof}
اسیں وکھاواں گے کہ کسے وی $\gamma \leq \alpha$ لئی، $\gamma$ تک دا
اکّو اک لگ بھگ روپ اے۔

پہلاں اسیں یکتائی ثابت کردے آں۔ منّ لو $g$ تے $h$ بالترتیب $\gamma$ تے
$\delta$ تک دے لگ بھگ روپ نیں۔ اوہناں دے مدخلاں اُتے لامتناہی درجیاں
والا استقرا وکھاندا اے کہ $g(\beta) = h(\beta)$ کسے وی
$\beta \in \dom{g} \cap \dom{h} = \gamma \cap \delta = \min(\gamma, \delta)$
لئی۔ سو ساڈے لگ بھگ روپ (جے موجود ہون) یکتا نیں تے ساریاں قدراں اُتے
اکّو نتیجہ دیندے نیں۔

موجودگی ثابت کرن لئی ہُن اسیں سادہ لامتناہی استقرا
(\olref[ordinals][opps]{simpletransrecursion}) ترتیبی عدداں
$\delta \leq \alpha$ اُتے ورتدے آں۔

صاف اے کہ خالی فنکشن $\emptyset$ تک دا لگ بھگ روپ اے۔

جے $g$، $\gamma$ تک دا لگ بھگ روپ اے، تاں $g \cup
\{\tuple{\gamma, \tau(g)}\}$، $\ordsucc{\gamma}$ تک دا لگ بھگ روپ اے۔

جے $\gamma$ حدّی ترتیبی عدد اے تے $g_\delta$، $\delta$ تک دا لگ بھگ
روپ اے سارے $\delta < \gamma$ لئی، تاں منّ لو
$g = \bigcup_{\delta \in \gamma} g_\delta$۔ ایہ فنکشن اے، کیوں جے ساڈے
وکھ وکھ $g_\delta$ ساریاں قدراں اُتے اکّو نتیجہ دیندے نیں۔ تے جے
$\delta \in \gamma$، تاں $g(\delta) = g_{\ordsucc{\delta}}(\delta) =
\tau(\funrestrictionto{g_{\ordsucc{\delta}}}{\delta}) =
\tau(\funrestrictionto{g}{\delta})$۔

ایس نال لامتناہی درجیاں اُتے استقرا والا ثبوت پورا ہویا۔
\end{proof}

جے اسیں فنکشن دی تھاں اک \emph{ٹرم} تعریف کرن دی اجازت دیئیے، تاں
پچھلے نتیجے توں حدّ $\alpha$ ہٹا سکدے آں۔ ہِیٹھلے نتیجے دے بیان تے ثبوت
وچ جدوں $\sigma$ اک ٹرم ہووے، اسیں
$\funrestrictionto{\sigma}{\alpha} = \Setabs{\tuple{\beta, \sigma(\beta)}}{\beta \in \alpha}$
رکھدے آں۔

\begin{thm}[عام بازگشت]\ollabel{transrecursionschema}
کسے وی ٹرم $\tau(x)$ لئی اسیں اک ٹرم $\sigma(x)$ صاف طور اُتے تعریف
کر سکدے آں، جیہدے لئی $\sigma(\alpha) = \tau(\funrestrictionto{\sigma}{\alpha})$
کسے وی ترتیبی عدد~$\alpha$ لئی۔
\end{thm}

\begin{proof}
ہر $\alpha$ لئی، \olref{transrecursionfun} نال
$\alpha$ تک دا اکّو اک لگ بھگ روپ، $f_\alpha$، اے،~$\tau$ لئی۔ $\sigma(\alpha)$ نوں
$f_{\ordsucc{\alpha}}(\alpha)$ تعریف کرو۔ ہُن:
	\begin{align*}
		\sigma(\alpha) &=
		f_{\ordsucc{\alpha}}(\alpha) \\&=
		\tau(\funrestrictionto{f_{\ordsucc{\alpha}}}{\alpha}) \\&=
		\tau(\Setabs{\tuple{\beta, f_{\ordsucc{\alpha}}(\beta)}}{\beta \in \alpha}) \\&=
		\tau(\Setabs{\tuple{\beta, f_{\ordsucc{\beta}}(\beta)}}{\beta \in \alpha})\\&=
		\tau(\funrestrictionto{\sigma}{\alpha})
	\end{align*}
دھیان دیو کہ $f_{\ordsucc{\beta}}(\beta) = f_{\ordsucc{\alpha}}(\beta)$
سارے $\beta < \alpha$ لئی، جیویں \olref{transrecursionfun} وچ۔
\end{proof}
\noindent
دھیان دیو کہ \olref{transrecursionschema} اک \emph{خاکہ} اے۔ ضروری گل
ایہ اے کہ اسیں $\sigma$ توں اک فنکشن، یعنی اک خاص قسم دا \emph{سیٹ}،
تعریف کرن دی اُمید نہیں رکھ سکدے، کیوں جے فیر $\dom{\sigma}$ سارے
ترتیبی عدداں دا سیٹ ہووے گا، جیہڑا بورالی فورتی دے پیراڈوکس
(\olref[ordinals][basic]{buraliforti}) دے اُلٹ اے۔

پر ہجے وکھاؤنا اے کہ \olref{transrecursionschema}، ساڈی $V_\alpha$ والے
سیٹاں دی تعریف دا جواز دیندا اے۔ ایہ جھٹ صاف نہ ہووے، پر لامتناہی
درجیاں والی بازگشت دے اک آخری، سادے روپ نال صاف ہو جاوے گا۔

\begin{thm}[سادہ بازگشت]\ollabel{simplerecursionschema}
کسے وی ٹرماں $\tau(x)$ تے $\theta(x)$ تے کسے وی سیٹ $A$ لئی، اسیں
اک ٹرم $\sigma(x)$ صاف طور اُتے تعریف کر سکدے آں جیہدے لئی:
\begin{align*}
	\sigma(\emptyset) &= A\\
	\sigma(\ordsucc{\alpha}) &= \tau(\sigma(\alpha)) &&
		\text{کسے وی ترتیبی عدد لئی }\alpha\\
	\sigma(\alpha) &= \theta(\ran{\funrestrictionto{\sigma}{\alpha}})&&
	\text{جدوں }\alpha\text{ حدّی ترتیبی عدد ہووے}
\end{align*}
\end{thm}

\begin{proof}
اسیں ٹرم $\xi(x)$ دی ایویں تعریف توں شروع کردے آں:
%t $\xi(x) = A$ if $x$ is not a function whose domain is an ordinal;
%otherwise let:
\[
	\xi(x) =
	\begin{cases}
			A &\text{جے $x$ ایہو جیہا فنکشن نہیں جیہدا}\\
			&\hspace{1em}\text{دائرہ ترتیبی عدد ہووے۔ نہیں تاں:}\\
			\tau(x(\alpha)) & \text{جے $\dom{x} = \ordsucc{\alpha}$}\\
			\theta(\ran{x}) & \text{جے $\dom{x}$ حدّی ترتیبی عدد اے}
	\end{cases}
\]
\olref{transrecursionschema} نال اک ٹرم $\sigma(x)$ اے جیہدے لئی
$\sigma(\alpha) = \xi(\funrestrictionto{\sigma}{\alpha})$ ہر ترتیبی
عدد $\alpha$ لئی۔ نالے $\funrestrictionto{\sigma}{\alpha}$ اک فنکشن
اے جیہدا دائرہ $\alpha$ اے۔ اسیں وکھاندے آں کہ $\sigma$ دیاں لوڑیندیاں خاصیتاں نیں،
سادہ لامتناہی استقرا (\olref[ordinals][opps]{simpletransrecursion}) نال۔

پہلاں، $\sigma(\emptyset) = \xi(\emptyset) = A$۔

فیر، $\sigma(\ordsucc{\alpha}) = \xi(\funrestrictionto{\sigma}{\ordsucc{\alpha}}) = \tau(\funrestrictionto{\sigma}{\ordsucc{\alpha}}(\alpha)) = \tau(\sigma(\alpha))$۔

آخر وچ، $\sigma(\alpha) = \xi(\funrestrictionto{\sigma}{\alpha}) = \theta(\ran{\funrestrictionto{\sigma}{\alpha}})$، جدوں $\alpha$ حدّی ترتیبی عدد اے۔
\end{proof}
\noindent
ہُن \olref[valpha]{defValphas} دا جواز دین لئی بس $A = \emptyset$،
$\tau(x) = \Pow{x}$ تے $\theta(x) = \bigcup x$ رکھو۔ آخر ایس نال
$V_\alpha$ والے سیٹاں دی تعریف دا جواز مل گیا!{}

\end{document}
